\documentclass[11pt]{article}
\usepackage[utf8]{inputenc}
\usepackage{geometry}
\usepackage{amsthm,marginnote}
\usepackage{microtype,amsmath,amssymb,verbatim,newpxtext}
\usepackage[euler-digits,euler-hat-accent]{eulervm}
\usepackage[mathscr]{euscript}

\usepackage{enumitem, tikz-cd,xr-hyper}
\usepackage[colorlinks=true]{hyperref}
\usepackage[noabbrev]{cleveref}

\usepackage{xparse}
\DeclareDocumentCommand{\dref}{om}{\IfNoValueTF{#1}{\cref{#2}}{#1, \cref{#1-#2}}}


\usepackage[style=alphabetic]{biblatex}
\bibliography{references}
\newcommand{\stackstaglink}[1]{\href{http://stacks.math.columbia.edu/tag/#1}{#1}}
\newcommand{\stackstag}[1]{\cite[\stackstaglink{#1}]{stacks}}
% find a way for making multiple tags

\newcommand{\reference}[1]{\marginnote{\footnotesize\textnormal{#1}}}

\date{}

% Numbered elements
\numberwithin{equation}{section}
\renewcommand{\thesubsection}{\thesection.\Alph{subsection}}

\newtheorem{theorem}[equation]{Theorem}
\newtheorem{lemma}[equation]{Lemma}
\newtheorem{proposition}[equation]{Proposition}
\newtheorem{corollary}[equation]{Corollary}

\theoremstyle{definition}
\newtheorem{definition}[equation]{Definition}
\newtheorem{history}[equation]{History}
\newtheorem{example}[equation]{Example}
\newtheorem{remark}[equation]{Remark}
\newtheorem{notation}[equation]{Notation}

\crefname{condition}{condition}{conditions}

% Commands
\newcommand{\st}{\; : \;}

\newcommand{\Z}{\mathbb{Z}}
\newcommand{\C}{\mathbb{C}}
\newcommand{\R}{\mathbb{R}}
\newcommand{\Q}{\mathbb{Q}}

\renewcommand{\O}{\mathscr{O}}

\DeclareMathOperator*\colim{colim}
\DeclareMathOperator{\coker}{coker}

\DeclareMathOperator{\supp}{supp}
\DeclareMathOperator{\rank}{rank}
\DeclareMathOperator{\height}{ht}

\DeclareMathOperator{\Convex}{Cnv}
\DeclareMathOperator{\Spec}{Spec}
\DeclareMathOperator{\Spa}{Spa}
\DeclareMathOperator{\Spv}{Spv}
\DeclareMathOperator{\Hom}{Hom}
\DeclareMathOperator{\End}{End}
\DeclareMathOperator{\Frac}{Frac}

% External documents
\externaldocument[ordered-groups-]{ordered-groups}
\externaldocument[valuations-]{valuations}

\title{Totally ordered abelian groups}

\begin{document}
\maketitle

\phantomsection
\label{section-phantom}
\tableofcontents

\section{Ordered sets}

\begin{definition}
Suppose $S$ is a set. An \emph{order} on $S$ is a relation $\geq$ on $S$ satisfying the following conditions.  
\begin{enumerate}[label=(\arabic{*})]
	\item (Reflexive) For any $x \in S$ we have $x \geq x$. 
	\item (Transitive) If $x \geq y$ and $y \geq z$, then $x \geq z$.
	\item (Antisymmetric) If $x \geq y$ and $y \geq x$, then $x = y$.  
\end{enumerate}
An order is \emph{total} if, for any $x$ and $y$ in $S$, either $x \geq y$ or $y \geq x$. 
\end{definition}

\begin{definition}
\begin{enumerate}[label=(\alph{*})]
	\item A \emph{ordered set} is a set equipped with an order. 
	\item A \emph{totally ordered set} is a set equipped with a total order. 
\end{enumerate}
\end{definition}

\begin{definition}
Suppose $S$ and $T$ are ordered sets and $\phi : S \to T$ is a function.
\begin{enumerate}[label=(\alph{*})]
\item $\phi$ is \emph{order preserving} if $x \geq y$ in $S$ implies $\phi(x) \geq \phi(y)$ in $T$. 
\item $\phi$ is \emph{order reversing} if $x \geq y$ in $S$ implies $\phi(x) \leq \phi(y)$ in $T$. 
\item $\phi$ is an \emph{order embedding} if $x \geq y$ in $S$ if and only if $\phi(x) \geq \phi(y)$ in $T$.
\item $\phi$ is an \emph{order isomorphism} if it is both a bijection and an order embedding. 
\end{enumerate}
\end{definition}

\begin{lemma}
\label{result:order-embeddings-injective}
Suppose $S$ and $T$ are ordered sets and $\phi : S \to T$ is an order embedding. Then $\phi$ is injective. 
\end{lemma}

\begin{proof}
Suppose $\phi(x) = \phi(y)$. Then we have $\phi(x) \geq \phi(y)$ and $\phi(y) \geq \phi(x)$, so $\phi$ being an order embedding implies that $x \geq y$ and $y \geq x$. Antisymmetry in $S$ then forces $x = y$. 
\end{proof}

\begin{lemma} \label{result:order-embedding-total-order}
Suppose $S$ and $T$ are ordered sets and $\phi : S \to T$ is a function. If $S$ is totally ordered, then the following are equivalent. 
\begin{enumerate}[label=(\alph{*})]
\item $\phi$ is an order embedding.
\item $\phi$ is injective and order preserving. 
\end{enumerate}
\end{lemma}

\begin{proof}
It follows from \cref{result:order-embeddings-injective} and definitions that (a) implies (b). Conversely, suppose $\phi$ is injective and order preserving. Suppose we have $x, y \in S$ such that $\phi(x) \geq \phi(y)$. We want to show that $x \geq y$. Since $S$ is totally ordered, we can assume for a contradiction that $x \lneq y$. Since $\phi$ is order preserving, this means that $\phi(x) \leq \phi(y)$. Antisymmetry in $T$ forces $\phi(x) = \phi(y)$, and $\phi$ being injective then implies $x = y$, contradicting our assumption that $x \lneq y$. 
\end{proof}

\begin{definition}
\label{definition:interval}
Suppose $S$ is an ordered set. For $a, b \in S$, we define the \emph{interval} $[a, b]$ by
\[ [a, b] = \{ x \in S : a \leq x \leq b \text{ or } b \leq x \leq a \}. \]
\end{definition}

\begin{definition}
\label{definition:convex-subset}
Let $S$ be an ordered set. A subset $T \subseteq S$ is \emph{convex} if, whenever $a, b \in T$, then $[a, b] \subseteq T$. 
\end{definition}

\begin{lemma}
\label{result:order-sum-order}
Let $S_1$ and $S_2$ be ordered sets. Then there exists on the disjoint union $S_1 \sqcup S_2$ a unique order satisfying both of the following conditions.
\begin{enumerate}[label=(\arabic{*})]
\item The inclusions of $S_1$ and $S_2$ are both order embeddings. 
\item All elements of $S_1$ are less than all elements of $S_2$. 
\end{enumerate}
\end{lemma}

\begin{proof}
Omitted.  
\end{proof}

\begin{definition}
\label{definition:ordered-sum}
If $S_1$ and $S_2$ are ordered sets, their \emph{ordered sum} $S_1 + S_2$ is the disjoint union $S \sqcup S_2$ equipped with the unique order from \cref{result:order-sum-order}.  
\end{definition}

\begin{lemma}
If $S_1$ and $S_2$ are totally ordered sets, then so is $S_1 + S_2$. 
\end{lemma}

\begin{proof}
Omitted.
\end{proof}

\section{Totally ordered abelian groups}

\begin{definition}
\label{definition:translation-invariant-relation}

Suppose $\Gamma$ is an abelian group. A relation $R \subseteq \Gamma \times \Gamma$ is \emph{translation  invariant} if, whenever $(a,b) \in R$, then
\[ (a + x, b+x) \in R  \]
for all $x \in \Gamma$. 
\end{definition}

\begin{lemma}
\label{result:orders-and-submonoids}

Let $\Gamma$ be an abelian group. There is an injective map
\[ \begin{tikzcd} \left\{ \text{translation invariant total orders on } \Gamma \right\} \ar{r} & \left\{ \text{submonoids of } \Gamma \right\} \end{tikzcd} \]
which carries a translation invariant total order $\geq$ on $\Gamma$ to the submonoid $\Gamma_{\geq  0} = \{x \in \Gamma : x \geq 0 \}$. Its image consists of the submonoids $S \subseteq \Gamma$ satisfying both of the following conditions.
\begin{enumerate}[label=(\arabic{*})]
\item \label[condition]{enumi:antisymmetry} If $x \in S$ and $-x \in S$, then $x = 0$. 
\item \label[condition]{enumi:totality} For every $x \in \Gamma$, either $x \in S$ or $-x \in S$. 
\end{enumerate}
\end{lemma}

\begin{proof}
Suppose we have a translation invariant total order $\geq$ on $\Gamma$ and $a, b \in \Gamma_{\geq 0}$. Since $a \geq 0$, translation-invariance guarantees that 
\[ a+b \geq b \geq 0. \]
Thus $\Gamma_{\geq 0}$ is a submonoid of $\Gamma$ and the map $\geq\,\mapsto \Gamma_{\geq 0}$ is well-defined. 
	
Suppose $\geq$ and $\geq'$ are both translation invariant total orders on $\Gamma$ such that $\Gamma_{\geq 0} = \Gamma_{\geq' 0}$. Suppose we have $a, b \in \Gamma$ such that $a \geq b$. By translation invariance of $\geq$, we have $a - b \geq 0$, so $a - b \in \Gamma_{\geq 0} = \Gamma_{\geq' 0}$. In other words, $a - b \geq' 0$, so by translation invariance of $\geq'$, have $a \geq' b$. Thus $a \geq b$ implies $a \geq' b$, and the converse implication is entirely symmetric. This shows that the map $\geq\,\mapsto \Gamma_{\geq 0}$ is injective.
	
Suppose we are given a submonoid $S \subseteq \Gamma$. We define a relation $\geq$ on $\Gamma$ by declaring $x \geq y$ if and only if $x - y \in S$. The fact that $S$ is a submonoid implies automatically that $\geq$ is reflexive, transitive, and translation invariant. It follows that $\geq$ is antisymmetric if and only if $S$ satisfies \cref{enumi:antisymmetry}, and that $\geq$ is total if and only if $S$ satisfies  \cref{enumi:totality} holds.
\end{proof}

\begin{lemma}
\label{result:positive-cones-disjoint}
Let $\Gamma$ be an abelian group and suppose $\geq$ and $\geq'$ are two total orders on $\Gamma$. If $\Gamma_{\geq 0} \subseteq \Gamma_{\geq' 0}$, then $\geq$ and $\geq'$ are equal. 
\end{lemma}

\begin{proof}
By the injecitivity assertion of \cref{result:orders-and-submonoids}, it is sufficient to show that $\Gamma_{\geq 0} = \Gamma_{\geq' 0}$. Suppose for a contradiction that $x \in \Gamma_{\geq' 0} \setminus \Gamma_{\geq 0}$. Then $-x \in \Gamma_{\geq 0} \subseteq \Gamma_{\geq' 0}$, so $x$ and $-x$ are both contained in $\Gamma_{\geq' 0}$. This means that $x = 0$, which contradicts our assumption that $x \notin \Gamma_{\geq 0}$. 
\end{proof}

\begin{definition}
\label{definition:totally-ordered-abelian-group}
A \emph{totally ordered abelian group} is an abelian group $\Gamma$ equipped with a translation invariant total order.
\end{definition}

\begin{lemma}
	\label{result:ordered-group-torsion-free}
	Suppose $\Gamma$ is a totally ordered abelian group. Then $\Gamma$ is torsion free.
\end{lemma}

\begin{proof}
	Suppose $x \in \Gamma$ and $x \gneq 0$. By translation invariance, we see that
	\[ \dotsb \gneq 3x \gneq 2x \gneq x \gneq 0. \]
	Indeed, if there were some positive integer $n$ such that $(n+1)x = nx$, then by subtracting $nx$ from both sides we would find that $x = 0$. In particular, we see that $x$ cannot be torsion. Similarly, if $x \lneq 0$, then 
	\[ 0 \gneq x \gneq 2x \gneq 3x \gneq \dotsb, \]
	and again $x$ cannot be torsion. 
\end{proof}

\begin{lemma}
\label{result:order-preserving-submonoid}
Suppose $\Gamma$ and $\Gamma'$ are totally ordered abelian groups and $\phi : \Gamma \to \Gamma'$ is a homomorphism of abelian groups. Then the following are equivalent.
\begin{enumerate}[label=(\alph{*})]
\item $\phi$ is order preserving.
\item $\phi(\Gamma_{\geq 0}) \subseteq \Gamma'_{\geq 0}$. 
\end{enumerate}
\end{lemma}

\begin{proof}
If $x \geq y$ in $\Gamma$, then $x - y \in \Gamma_{\geq 0}$, so $\phi(x-y) = \phi(x)-\phi(y) \in \phi(\Gamma_{\geq 0})$. The statement follows immediately.
\end{proof}

\begin{lemma}
\label{result:order-embedding-submonoid}
Suppose $\Gamma$ and $\Gamma'$ are totally ordered abelian groups and $\phi : \Gamma \to \Gamma'$ is a homomorphism of abelian groups. Then the following are equivalent.
\begin{enumerate}[label=(\alph{*})]
\item $\phi$ is an order embedding.
\item $\Gamma_{\geq 0} = \phi^{-1}(\Gamma'_{\geq 0})$. 
\end{enumerate}
\end{lemma}

\begin{proof}
Suppose $\phi$ is an order embedding. We know from \cref{result:order-preserving-submonoid} that $\Gamma_{\geq 0} \subseteq \phi^{-1}(\Gamma'_{\geq 0})$. For the reverse inclusion, suppose $x \in \phi^{-1}(\Gamma'_{\geq 0})$. In other words, we have $\phi(x) \geq 0$. We would like to show that $x \geq 0$, so suppose for a contradiction that $x \lneq 0$. Since $\phi$ is order preserving, this implies that $\phi(x) \leq \phi(0) = 0$, so antisymmetry forces $\phi(x) = 0$.  But $\phi$ is injective by \cref{result:order-embeddings-injective}, so $x = 0$, which contradicts our assumption that $x \lneq 0$.

Conversely, suppose $\Gamma_{\geq 0} = \phi^{-1}(\Gamma'_{\geq 0})$. It follows from \cref{result:order-preserving-submonoid} that $\phi$ is order preserving. Suppose $x, y \in \Gamma$ are such that $\phi(x) \geq \phi(y)$. Then $\phi(x) - \phi(y) = \phi(x-y) \in \Gamma_{\geq 0}'$, or, in other words, $x - y \in \phi^{-1}(\Gamma_{\geq 0}')$. Since $\phi^{-1}(\Gamma'_{\geq 0}) = \Gamma_{\geq 0}$, we conclude that $x \geq y$. 
\end{proof}

\section{Constructions}

\begin{example}
The real numbers $\R$, regarded as an abelian group under addition and endowed with the standard order, form a totally ordered abelian group. 
\end{example}

\begin{example}
Regard $\R[\epsilon] = \{ a_0 + a_1 \epsilon + \dotsb + a_n \epsilon^n : a_0, \dotsc, a_n \in \R \}$ as an abelian group under addition. Then the submonoid
\[ S = \{ a_k\epsilon^k + a_{k+1} \epsilon + \dotsb + a_n \epsilon^n : a_k \gneq 0 \} \]
defines a translation invariant total order on $\R[\epsilon]$ under the correspondence of \cref{result:orders-and-submonoids}. In the resulting totally ordered abelian group $\R[\epsilon]$, the element $\epsilon$ ``behaves like a positive infinitesimal.'' For example, we have 
\[ 0 \lneq \dotsb \lneq \epsilon^2 \lneq \epsilon \lneq a \]
for any $a \in \R$. 
\end{example}

\begin{example}
\label{example:subgroup-order}
If $\Gamma$ is a totally ordered abelian group and $\Gamma' \subseteq \Gamma$ is any subgroup, then $\Gamma'$ is a totally ordered abelian group with the induced total order. 
\end{example}

\begin{example}
\label{example:divisible-closure}
Suppose $\Gamma$ is a totally ordered abelian group. It is torsion free by \cref{result:ordered-group-torsion-free}, and there exists on the divisible closure $\Q \otimes \Gamma$ a unique translation invariant total order such that $\Gamma \hookrightarrow \Q \otimes \Gamma$ is an order embedding. It satisfies 
\[ (\Q \otimes \Gamma)_{\geq 0} = \{ x \in \Q \otimes \Gamma \st nx \in \Gamma_{\geq 0} \text{ for some } n \geq 1 \}. \]
\end{example}

\begin{proof}[Proof of existence]
We use \cref{result:orders-and-submonoids}. Observe that the set 
\[ S = \{x \in \Q \otimes \Gamma \st nx \in \Gamma_{\geq 0} \text{ for some } n \geq 1 \} \] is a submonoid of $\Q \otimes \Gamma$. Indeed, it clearly contains zero, and if we have $x, y \in S$, then there exist $m, n \geq 1$ such that $mx, ny \in \Gamma_{\geq 0}$. Then
\[ mn(x+y) = n(mx) + m(ny) \in \Gamma_{\geq 0} \]
showing that $x+y \in S$. If $x \in S$ and $-x \in S$, then there exist $m, n \geq 1$ such that $mx \in \Gamma_{\geq 0}$ and $-nx \in \Gamma_{\geq 0}$. Then $mnx$ and $-mnx$ are both in $\Gamma_{\geq 0}$, which means that $mnx = 0$. Since $\Q \otimes \Gamma$ is clearly torsion free, this means that $x  = 0$, proving that \cref{enumi:antisymmetry} is satisfied. Finally, to check that \cref{enumi:totality} is satisfied, suppose that we have an arbitrary element $x \in \Q \otimes \Gamma$. Then there exists some integer $n \geq 1$ such that $nx \in \Gamma$. Then either $nx \in \Gamma_{\geq 0}$ or $-nx \in \Gamma_{\geq 0}$.  In the first case, we have $x \in S$, and in the second, we have $-x \in S$. Thus there exists a unique translation invariant total order $\geq$ on $\Q \otimes \Gamma$ such that $(\Q \otimes \Gamma)_{\geq 0} = S$. 
\end{proof}

\begin{proof}[Proof of uniqueness]
Let $S$ be the submonoid described in the proof of existence above, and suppose $\geq$ is an arbitrary translation invariant total order of $\Q \otimes \Gamma$ such that $\Gamma \hookrightarrow \Q \otimes \Gamma$ is an order embedding. Then $(\Q \otimes \Gamma)_{\geq 0} \cap \Gamma = \Gamma_{\geq 0}$ by \cref{result:order-embedding-submonoid}. We claim that this implies that $S = (\Q \otimes \Gamma)_{\geq 0}$, and then the injectivity assertion of \cref{result:orders-and-submonoids} completes the proof of uniqueness. 

To see that $S = (\Q \otimes \Gamma)_{\geq 0}$, suppose first that $x \in S$. Then $nx \in \Gamma_{\geq 0} = (\Q \otimes \Gamma)_{\geq 0} \cap \Gamma$ for some $n \geq 1$. In particular, $nx \in (\Q \otimes \Gamma)_{\geq 0}$, so clearly we must have $x \in (\Q \otimes \Gamma)_{\geq 0}$. On the other hand, if $x \in (\Q \otimes \Gamma)_{\geq 0}$, there exists some integer $n \geq 1$ such that $nx \in \Gamma$. Then $nx \in (\Q \otimes \Gamma)_{\geq 0} \cap \Gamma = \Gamma_{\geq 0}$, so $x \in S$ by definition of $S$.
\end{proof}

\begin{example}
\label{example:hahn-product}
Suppose $I$ is a totally ordered index set and $(\Gamma_i)_{i \in I}$ a family of totally ordered abelian groups. For any $\gamma = (\gamma_i)_{i \in I} \in \prod_{i \in I} \Gamma_i$, we define 
\[ \supp(\gamma) = \{ i \in I : \gamma_i \neq 0 \}. \]
The \emph{Hahn product} of the family $(\Gamma_i)_{i \in I}$ is defined by
\[ \prod_{i \in I}' \Gamma_i = \left\{ \gamma \in \prod_{i \in I} \Gamma_i \st \supp(\gamma) \text{ is well ordered} \right\}. \]
This is a subgroup of $\prod_{i \in I} \Gamma_i$, and is endowed with a total order $\geq$ corresponding (under \cref{result:orders-and-submonoids}) to the submonoid
\begin{equation} \label{eqn:hahn-product-order-cone} \left\{ \gamma = (\gamma_i)_{i \in I} \in \prod_{i \in I}' \Gamma_i \st \gamma = 0 \text{ or } \gamma_{j} \geq 0 \text{ for } j = \min \supp(\gamma)  \right\}.
\end{equation}
This total order $\geq$ is called the \emph{lexicographic order} on the Hahn product.
\end{example}

\begin{proof}[Proof that the Hahn product is a subgroup]
It is clear that the Hahn product contains the zero element. Suppose we have $\gamma = (\gamma_i)_{i \in I}$ and $\gamma' = (\gamma'_i)_{i \in I}$ such that $\supp(\gamma)$ and $\supp(\gamma')$ are well ordered. Let $S$ be a nonempty subset of $\supp(\gamma + \gamma')$, and suppose for a contradiction that it has no minimal element. Then we can find an infinite, strictly decreasing sequence $i_0 \gneq i_1 \gneq \dotsb$ in $S$. In other words, we have $\gamma_{i_k} + \gamma'_{i_k} \neq 0$ for all $k = 0, 1, \dotsc$. 
	
Since $\supp(\gamma)$ is well-ordered, no infinite subsequence of $i_0 \gneq i_1 \gneq \dotsb$ can be contained in $\supp(\gamma)$, which means that there exists some $K$ sufficiently large such that $\gamma_{i_k} = 0$ for all $k \geq K$. For the same reasons, we can make $K$ larger if necessary in order to have that $\gamma'_{i_k} = 0$ for all $k \geq K$ as well. But this clearly forces $\gamma_{i_k} + \gamma'_{i_k} = 0$, which is a contradiction.  
\end{proof}

\begin{proof}[Proof that the lexicographic order is a total order on the Hahn product]
We need to verify that the submonoid in \cref{eqn:hahn-product-order-cone} satisfies the conditions of \cref{result:orders-and-submonoids}. \Cref{enumi:totality} is clear. For \cref{enumi:antisymmetry}, suppose both $\gamma = (\gamma_{i})_{i \in I}$ and its negative are in the submonoid of \cref{eqn:hahn-product-order-cone}. Suppose there exists $j \in \supp(\gamma)$. Since $\supp(\gamma) = \supp(-\gamma)$, we see that $\gamma_j \geq 0$ and $-\gamma_j \geq 0$, so $\gamma_j = 0$. This contradicts our choice of $j$, so in fact $\supp(\gamma) = \emptyset$. In other words, $\gamma = 0$. 
\end{proof}

\section{Convex subgroups}

\begin{definition}
\label{definition:convex-subgroup}
Suppose $\Gamma$ is a totally ordered abelian group. A \emph{convex subgroup} (or an \emph{isolated subgroup}) of $\Gamma$ is a subgroup which is convex in the sense of \cref{definition:convex-subset}.
\end{definition}

\begin{lemma}
	\label{result:convexity-positive-sufficient}
	Suppose $\Gamma$ is a totally ordered abelian group and $\Delta$ is a subgroup. Then the following are equivalent.
	\begin{enumerate}[label=(\alph{*})]
		\item $\Delta$ is convex.
		\item Whenever $a \geq x \geq 0$ and $a \in \Delta$, then $x \in \Delta$ also. 
	\end{enumerate}
\end{lemma}

\begin{proof}
	Suppose we have $a \geq x \geq b$ where $a, b \in \Delta$. By translation invariance, we have 
	\[ a - b \geq x - b \geq 0 \]
	and $a - b \in \Delta$, so the condition in (b) implies that $x - b \in \Delta$. But since $b \in \Delta$, this means that $x \in \Delta$ as well. 
\end{proof}

\begin{lemma}
\label{result:kernels-convex}
Suppose $\phi : \Gamma \to \Gamma'$ is an order preserving homomorophism of totally ordered abelian groups. Then $\ker(\phi)$ is convex. 
\end{lemma}

\begin{proof}
We use the criterion of \cref{result:convexity-positive-sufficient}. Suppose we have $a \geq x \geq 0$ in $\Gamma$ with $a \in \ker(\phi)$. Then \[ 0 = \phi(a) \geq \phi(x) \geq 0,  \] so antisymmetry forces $\phi(x) = 0$. 
\end{proof}

\begin{lemma}
\label{result:convex-inside-convex}
Suppose $\Gamma$ is a totally ordered abelian group, $\Delta$ is a convex subgroup of $\Gamma$, and $\Delta'$ is a convex subgrup of $\Delta$. Then $\Delta'$ is a convex subgroup of $\Gamma$. 
\end{lemma}

\begin{proof}
We use the criterion of \cref{result:convexity-positive-sufficient}. Suppose we have $a \geq x \geq 0$ in $\Gamma$ where $a \in \Delta'$. Since $\Delta' \subseteq \Delta$ and $\Delta$ is a convex subgroup of $\Delta$, we see that $x \in \Delta$. Then since $\Delta'$ is convex in $\Delta$, we have $x \in \Delta'$ also. 
\end{proof}

\begin{lemma}
	\label{result:convex-subgroups-totally-ordered}
	Let $\Gamma$ be a totally ordered abelian group and suppose $\Delta$ and $\Delta'$ are two convex subgroups of $\Gamma$. Then either $\Delta \subseteq \Delta'$ or $\Delta' \subseteq \Delta$. 
\end{lemma}

\begin{proof}
	Suppose $\Delta \not\subseteq \Delta'$. Then there exists $a \in \Delta \setminus \Delta'$, and by replacing $a$ with $-a$ if necessary, we can assume $a \geq 0$. Suppose $x \in \Delta'$. We claim that $x \in \Delta$. By replacing $x$ with $-x$ if necessary, we can assume that $x \geq 0$. If $x \geq a$, then we would have $x \geq a \geq 0$ and convexity of $\Delta'$ would force $a \in \Delta'$, contradicting our choice of $a$. Thus we must have $a \geq x \geq 0$, and convexity of $\Delta$ now implies that $x \in \Delta$. 
\end{proof}


\begin{lemma}
	\label{result:intersection-convex}
	Suppose $\Gamma$ is a totally ordered abelian group and $(\Delta_i)_{i \in I}$ is a family of convex subgroups of $\Gamma$. Then $\Delta = \bigcap_{i \in I} \Delta_i$ is a convex subgroup of $\Gamma$. 
\end{lemma}

\begin{proof}
	We use the criterion of \cref{result:convexity-positive-sufficient}. Suppose we have $a \geq x \geq 0$ in $\Gamma$ where $a \in \Delta$. Then $a \in \Delta_i$ for all $i$, so $x \in \Delta_i$ for all $i$, so $x \in \Delta$. 
\end{proof}

\begin{lemma}
	\label{result:union-convex}
	Suppose $\Gamma$ is a totally ordered abelian group and $(\Delta_i)_{i \in I}$ is a nonempty family of convex subgroups of $\Gamma$. Then $\Delta = \bigcup_{i \in I} \Delta_i$ is a convex subgroup of $\Gamma$. 
\end{lemma}

\begin{proof}
Clearly $0 \in \Delta$. If $a, b \in \Delta$, then there exist $i, j \in I$ such that $a \in \Delta_i$ and $b \in \Delta_j$. By \cref{result:convex-subgroups-totally-ordered}), we can assume without loss of generality that $\Delta_i \subseteq \Delta_j$. Then $a - b \in \Delta_j  \subseteq \Delta$ also, proving that $\Delta$ is a subgroup.  
To see that it is convex, we use the criterion of \cref{result:convexity-positive-sufficient}. Suppose we have $a \geq x \geq 0$ in $\Gamma$ where $a \in \Delta$. Then there exists $i$ such that $a \in \Delta_i$, so then $x \in \Delta_i \subseteq \Delta$ also. 
\end{proof}

\begin{lemma}
\label{result:convex-subgroup-divisible}
Let $\Gamma$ be a totally ordered abelian group and $\Delta$ a convex subgroup. If $\Gamma$ is divisible, then so is $\Delta$. 
\end{lemma}

\begin{proof}
Suppose that $x \in \Gamma$ is some element such that $nx \in \Delta$ for some $n \geq 1$. We want to show that $x \in \Delta$. We can assume without loss of generality that $x \geq 0$. Translation invariance implies that
\[ \dotsb \geq nx \geq (n-1)x \geq \dotsb \geq 2x \geq x \geq 0, \]
so $x \in [0, nx]$. Convexity of $\Delta$ implies that $x \in \Delta$. 
\end{proof}

\begin{definition}
Let $\Gamma$ be a totally ordered abelian group. Given any subset $S \subseteq \Gamma$, we define the \emph{convex subgroup of $\Gamma$ generated by $S$} to be the intersection of all convex subgroups of $\Gamma$ containing $S$. By \cref{result:intersection-convex}, this is the smallest convex subgroup of $\Gamma$ containing $S$. 
\end{definition}

\begin{lemma}
	\label{result:convex-subgroup-generated-characterization}
Suppose $\Gamma$ is a totally ordered abelian group, $H$ is a subgroup of $\Gamma$ and $\Delta$ is the convex subgroup of $\Gamma$ generated by $H$. Then 
\[ \Delta = \bigcup_{h \in H_{\geq 0}} [-h, h]. \]
\end{lemma}

\begin{proof}
It is clear that $S = \bigcup_{h \in H_{\geq 0}} [-h, h]$ must be contained in every convex subgroup containing $H$, so $S \subseteq \Delta$. To prove that $S = \Delta$, we need to show that $S$ is actually a convex subgroup. It clearly contains zero and is stable under negation. If we have $x, x' \in S$, then there exist $h, h' \in H_{\geq 0}$ such that $x \in [-h, h]$ and $x' \in [-h', h']$. Then
\[ h + h' \geq x + x' \geq -h - h'. \]
In other words, we have $x + x' \in [-(h+h'), h+h'] \subseteq S$, showing that $S$ is stable under addition as well. Finally, to see that $S$ is convex, we use the criterion of \cref{result:convexity-positive-sufficient}. Suppose we have $a \geq x \geq 0$ with $a \in S$. Then there exists some $h \in H_{\geq 0}$ such that $a \in [-h, h]$. But then clearly $x \in [-h, h]$ also, so $x \in S$ as well. 
\end{proof}

\section{Ordered set of convex subgroups}

\begin{definition}
\label{definition:set-of-convex}
Let $\Gamma$ be a totally ordered abelian group. We write $\Convex(\Gamma)$ for the set of all convex subgroups of $\Gamma$. By \cref{result:convex-subgroups-totally-ordered}, this set is totally ordered by inclusion; evidently $0$ and $\Gamma$ are the minimum and maximum elements of $\Convex(\Gamma)$, respectively. We write $\Convex^\circ(\Gamma)$ for $\Convex(\Gamma) \setminus \{\Gamma\}$. 
\end{definition}

\begin{lemma}
\label{result:pullback-convex-subgroups}
Suppose $\phi : \Gamma \to \Gamma'$ is an order preserving homomorphism between totally ordered abelian groups. Then $\Delta' \mapsto \phi^{-1}(\Delta')$ is an order preserving map $\phi^{-1} : \Convex(\Gamma') \to \Convex(\Gamma)$.
\end{lemma}

\begin{proof}
Suppose $\Delta'$ is a convex subgroup of $\Gamma'$ and suppose $a, b \in \phi^{-1}(\Delta')$ and that $a \leq x \leq b$ for some $x \in \Gamma$. Since $\phi$ is order preserving, we have $\phi(a) \leq \phi(x) \leq \phi(b)$ in $\Gamma'$. Since $\phi(a), \phi(b) \in \Delta'$ and $\Delta'$ is convex, this means that $\phi(x) \in \Delta'$, so $x \in \phi^{-1}(\Delta')$. In other words, $\phi^{-1}(\Delta')$ is a convex subgroup of $\Gamma$. The fact that $\Delta' \mapsto \phi^{-1}(\Delta')$ is inclusion preserving is clear. 
\end{proof}

\begin{lemma}
\label{result:convex-subgroups-divisible-closure}
Suppose $\Gamma$ is a totally ordered abelian group. Then the natural order embedding $\Gamma \hookrightarrow \Q \otimes \Gamma$ of \cref{example:divisible-closure} induces an order isomorphism $\Convex(\Q \otimes \Gamma) \to \Convex(\Gamma)$. Its inverse is given by $\Delta \mapsto \Q \otimes \Delta$. 
\end{lemma}

\begin{proof}
Since $\Convex(\Q \otimes \Gamma)$ is totally ordered by \cref{result:convex-subgroups-totally-ordered} and $\Convex(\Q \otimes \Gamma) \to \Convex(\Gamma)$ is inclusion preserving by \cref{result:pullback-convex-subgroups}, it is sufficient to prove that $\Convex(\Q \otimes \Gamma) \to \Convex(\Gamma)$ is bijective (cf. \cref{result:order-embedding-total-order}).  
	
For injectivity, suppose that $\Delta$ and $\Delta'$ are convex subgroups of $\Q \otimes \Gamma$ such that $\Delta \cap \Gamma = \Delta' \cap \Gamma$. Since $\Convex(\Q \otimes \Gamma)$ is totally ordered, we can assume without loss of generality $\Delta \subseteq \Delta'$. We will show that $\Delta = \Delta'$. Suppose $x \in \Delta'$. We can assume without loss of generality that $x \geq 0$. Then there exists some $n \geq 1$ such that $nx \in \Gamma$. Then $nx \in \Delta' \cap \Gamma = \Delta \cap \Gamma$. But $\Delta$ is divisible by \cref{result:convex-subgroup-divisible}, so $x \in \Delta$ also. 

For surjectivity, suppose $\Delta$ is a convex subgroup of $\Gamma$. Then $\Q \otimes \Delta$ is a convex subgroup of $\Q \otimes \Gamma$. To see this, suppose that $a \geq x \geq b$ in $\Q \otimes \Gamma$ where $a, b \in \Q \otimes \Delta$. There exists some $n \geq 1$ such that $na, nx, nb \in \Gamma$, and in fact we have $na, nb \in \Delta$. Translation invariance implies that $na \geq nx \geq nb$, so convexity of $\Delta$ tells us that $nx \in \Delta$. This means that $x \in \Q \otimes \Delta$.

Finally, note that $(\Q \otimes \Delta) \cap \Gamma = \Delta$. Indeed, clearly $\Delta \subseteq (\Q \otimes \Delta) \cap \Gamma$. For the reverse inclusion, suppose that $x \in (\Q \otimes \Delta) \cap \Gamma$ and assume without loss of generality that $x \geq 0$. Then $x \in \Gamma$ and there exists $n \geq 1$ such that $nx \in \Delta$. Translation invariance implies that
\[ \dotsb \geq nx \geq (n-1)x \geq \dotsb \geq x \geq 0 \]
so $x \in [0, nx]$. Convexity of $\Delta$ thus implies that $x \in \Delta$. 
\end{proof}

\section{Quotients by convex subgroups}

\begin{lemma}
\label{result:quotient-order}
Suppose $\Gamma$ is a totally ordered abelian group, $\Delta$ is a convex subgroup, and $\pi : \Gamma \to \Gamma/\Delta$ is the quotient map. Then there exists a unique translation invariant total order on $\Gamma/\Delta$ such that $\pi$ is order preserving. Moreover, we have 
\[ (\Gamma/\Delta)_{\geq 0} = \pi(\Gamma_{\geq 0}). \]
\end{lemma}

\begin{proof}
For existence, we will prove that the submonoid $\pi(\Gamma_{\geq 0})$ of $\Gamma/\Delta$ satisfies the conditions required to be in the image of the map in \cref{result:orders-and-submonoids}. \Cref{enumi:totality} is clear: every element of $\Gamma/\Delta$ is of the form $\pi(x)$ for some $x \in \Gamma$, and either $x \geq 0$, in which case $\pi(x) \in \pi(\Gamma_{\geq 0})$, or else $-x \geq 0$, in which case $-\pi(x) = \pi(-x) \in \pi(\Gamma_{\geq 0})$. By unwinding \cref{enumi:antisymmetry}, we find that we need to show the following: if $x \in \Gamma_{\geq 0}$ and there exists $y \in \Gamma_{\geq 0}$ such that $x + y \in \Delta$, then $x \in \Delta$ also. But notice that, by translation invariance, we have \[ x + y \geq x \geq 0 \] so convexity implies that $x \in \Delta$. 
	
For uniqueness, suppose $\geq'$ is an order on $\Gamma/\Delta$ which makes $\pi$ order preserving. We must have $\pi(\Gamma_{\geq 0}) \subseteq (\Gamma/\Delta)_{\geq' 0}$ by \cref{result:order-preserving-submonoid}. Then the existence proof above plus \cref{result:positive-cones-disjoint} guarantee that we must have $\pi(\Gamma_{\geq 0}) = (\Gamma/\Delta)_{\geq' 0}$. 
\end{proof}

\begin{lemma}
\label{result:quotient-order-positives}
Suppose $\Gamma$ is a totally ordered abelian group, $\Delta$ is a convex subgroup, and $\pi : \Gamma \to \Gamma/\Delta$ is the quotient map. If $x \in \Gamma$ is such that $\pi(x) \geq 0$, then either $x \geq 0$ or $x \in \Delta$. 
\end{lemma}

\begin{proof}
Suppose that $x \lneq 0$. Since $\pi$ is order preserving, we have $\pi(x) \leq 0$, so antisymmetry in $\Gamma/\Delta$ forces $\pi(x) = 0$. In other words, $x \in \ker(\pi) = \Delta$. 
\end{proof}

\begin{lemma}
\label{result:quotient-convex-subgroups}
Suppose $\Gamma$ is a totally ordered abelian group, $\Delta$ is a convex subgroup, and $\pi : \Gamma \to \Gamma/\Delta$ is the quotient map. Then $\pi^{-1} :  \Convex(\Gamma/\Delta) \to \Convex(\Gamma)$ is injective, with image consisting of all convex subgroups of $\Gamma$ which contain $\Delta$. 
\end{lemma}

\begin{proof}
Since $\pi^{-1}$ defines a bijection between subgroups of $\Gamma/\Delta$ and subgroups of $\Gamma$ containing $\Delta$, restricting $\pi^{-1}$ to $\Convex(\Gamma/\Delta)$ still yields an injective map whose image is contained in the set of subgroups containing $\Delta$. Moreover, applying $\pi^{-1}$ to a \emph{convex} subgroup yields another convex subgroup by \cref{result:pullback-convex-subgroups}.
	
To complete the proof, we need to show that, if $\Delta'$ is a convex subgroup of $\Gamma$ containing $\Delta$, then $\Delta'/\Delta$ is a convex subgroup of $\Gamma/\Delta$. We use the criterion of \cref{result:convexity-positive-sufficient}. Suppose we have $\pi(a) \geq \pi(x) \geq 0$ where $a \in \Delta' \cap \Gamma_{\geq 0}$ and $x \in \Gamma_{\geq 0}$. Then $\pi(a-x) \geq 0$. By \cref{result:quotient-order-positives}, there are two cases. First, it could be that $a - x \geq 0$. In this case we have $a \geq x \geq 0$, so convexity forces $x \in \Delta'$ and we are done. The second case is that $a - x \in \Delta$. But then $\pi(x) = \pi(a) \in \Delta'/\Delta$, so again we are done. 
\end{proof}

\section{Height of a totally ordered abelian group}

\begin{definition}
	\label{definition:height}
	The \emph{height} (or \emph{rank}) of $\Gamma$, denoted $\height(\Gamma)$, is the cardinality of the set $\Convex^\circ(\Gamma)$ of proper convex subgroups of $\Gamma$.
\end{definition}

\begin{lemma}
Suppose $\Gamma$ is a totally ordered abelian group. Then $\height(\Q \otimes \Gamma) = \height(\Gamma)$. 
\end{lemma}

\begin{proof}
This follows immediately from \cref{result:convex-subgroups-divisible-closure}. 
\end{proof}

\begin{lemma}
	Suppose $\Gamma$ is a totally ordered abelian group and $\Delta$ a convex subgroup. Then
	\[ \height(\Gamma) = \height(\Delta) + \height(\Gamma/\Delta). \]
\end{lemma}

\begin{proof}
	Let $\pi : \Gamma \to \Gamma/\Delta$ be the quotient map. The sum of two cardinals is the cardinality of the disjoint union, so we want to show that \[ \Convex^\circ(\Delta) \sqcup \Convex^\circ(\Gamma/\Delta) \] bijects with $\Convex^\circ(\Gamma)$. In fact, we will produce an order isomorphism
	\[ \begin{tikzcd} \Convex^\circ(\Delta) + \Convex(\Gamma/\Delta) \ar{r}{f} & \Convex(\Gamma) \end{tikzcd} \] 
	such that $f(\Gamma/\Delta) = \Gamma$, where the domain is the ordered sum of \cref{definition:ordered-sum}. 
	
	On $\Convex(\Gamma/\Delta)$, we define $f$ to be $\pi^{-1} : \Convex(\Gamma/\Delta) \to \Convex(\Gamma)$. This is order preserving from \cref{result:pullback-convex-subgroups}. Also, if $\Delta'$ is a proper convex subgroup of $\Delta$, it is also a convex subgroup of $\Gamma$ (cf. \cref{result:convex-inside-convex}), and we define $f(\Delta')$ to be $\Delta'$ regarded as a convex subgroup of $\Gamma$. This defines the function $f$, which is clearly order preserving and we have $f(\Gamma/\Delta) = \Gamma$. 
	
	On the other hand, consider the function 
	\[ \begin{tikzcd} \Convex(\Gamma) \ar{r}{g} & (\Convex(\Delta) \setminus \{\Delta\}) + \Convex(\Gamma/\Delta) \end{tikzcd} \] defined by
	\[ g(\Delta') = \begin{cases} \Delta'/\Delta \in \Convex(\Gamma/\Delta) & \text{if } \Delta' \supseteq \Delta \\ 
	\Delta' \in \Convex^\circ(\Delta) & \text{if } \Delta \supsetneq \Delta' \end{cases}  \]
	We know from \cref{result:quotient-convex-subgroups} that $\Delta'/\Delta$ is in fact a convex subgroup of $\Gamma/\Delta$ when $\Delta' \supseteq \Delta$, so $g$ is well defined. It is clear that $g$ is order preserving that it is an inverse for $f$. 
\end{proof}

\begin{lemma}
	\label{result:height-1-ordered-groups}
	\reference{\cite[VI.4, proposition 8]{bourbaki_ac} \\
	\medskip
	\cite[IV, theorem 1]{fuchs}}

	Suppose $\Gamma$ is a totally ordered abelian group. Then the following are equivalent.
	\begin{enumerate}[label=(\alph{*})]
		\item $\height(\Gamma) \leq 1$. 
		\item For all $a \gneq 0$ and $x \geq 0$ in $\Gamma$, there exists some $n \geq 1$ such that $na \geq x$. 
		\item There exists an order embedding $\Gamma \hookrightarrow \R$. 
	\end{enumerate}
\end{lemma}

\begin{proof}
	(a) $\implies$ (b). Fix $a \gneq 0$ and $x \geq 0$ in $\Gamma$. Let $\Delta_a$ be the convex subgroup of $\Gamma$ generated by $a$. By \cref{result:convex-subgroup-generated-characterization}, we know that
	\[ \Delta_a = \bigcup_{n \geq 1} [-na, na]. \]
	Then $\Delta_a$ is a nonzero convex subgroup of $\Gamma$, so, since $\height(\Gamma) \leq 1$, we must have $\Delta_a = \Gamma$. Thus $x \in \Delta_a$, which means precisely that there exists $n \geq 1$ such that $na \geq x$. 
	
	\medskip
	\noindent
	(b) $\implies$ (c). If $\Gamma = 0$, there is nothing to show, so let us assume that $\Gamma \neq 0$. Fix an element $a \gneq 0$ in $\Gamma$. For any element $x \in \Gamma$, we define the following. 
	\[ \begin{aligned} U(x) &= \{ m/n \in \Q \st m, n \in \Z, n \geq 1, ma \geq nx \} \\
	L(x) &= \{ m/n \in \Q \st m, n \in \Z, n \geq 1, ma \lneq nx \} \end{aligned} \]
	The pair $(L(x),U(x))$ is a Dedekind cut in $\Q$, and we set 
	\[ \phi(x) = \inf U(x) = \sup L(x). \]
	Then $\phi$ defines an injective and order preserving homomorphism $\Gamma \to \R$ with the property that $\phi(a) = 1$. Details omitted. By \cref{result:order-embedding-total-order}, $\phi$ must be an order embedding. 
	
	\medskip
	\noindent
	(c) $\implies$ (a). Observe that $\height(\R) = 1$. Indeed, suppose $\Delta$ is a nonzero convex subgroup of $\R$ and $a \gneq 0$ in $\Delta$. Let $\Delta_a$ be the convex sugbroup of $\Gamma$ generated by $a$. Then, by \cref{result:convex-subgroup-generated-characterization}, we have
	\[ \Delta_a = \bigcup_{n \geq 1} [-na, na] \subseteq \Delta. \]
	It is clear that every $x \in \R$ is contained in $[-na, na]$ for some $n \geq 1$, so $\R = \Delta_a$, forcing $\Delta = \R$. It follows that if there exists an order embedding $\Gamma \hookrightarrow \R$, then $\height(\Gamma) \leq \height(\R) = 1$. 
\end{proof}


\section{Categorical properties}

\begin{definition}
The \emph{category of totally ordered abelian groups} is the category whose objects are totally ordered abelian groups and whose morphisms are order preserving homomorphisms. 
\end{definition}

\begin{remark}
The set of order preserving homomorphisms between two totally ordered abelian groups is a submonoid of the group of all (not necessarily order preserving) homomorphisms between them: it contains the zero map and is stable under addition, but clearly the negative of an order preserving map is \emph{not} order preserving. 
\end{remark}

\begin{lemma}
\label{result:categorical-zero-kernels-cokernels}
The category of totally ordered abelian groups has a zero object, kernels, and cokernels. 
\end{lemma}

\begin{proof}
Clearly the trivial group is a zero object. Suppose $\phi : \Gamma \to \Gamma'$ is an order preserving homomorphism. Then the group-theoretic kernel $\ker(\phi)$ with the induced order (as in \cref{example:subgroup-order}) is the kernel of $\phi$ in the category of totally ordered abelian groups. 

Let $\Delta$ denote the convex subgroup of $\Gamma'$ generated by $\phi(\Gamma)$ (i.e., the intersection of all convex subgroups containing $\phi(\Gamma)$, which is again a convex subgroup by \cref{result:intersection-convex}). Then the quotient map $\pi : \Gamma' \to \Gamma'/\Delta$ is a cokernel for $\phi$ in the category of totally ordered abelian groups. Indeed, suppose $\phi' : \Gamma' \to \Gamma''$ is an order preserving homomorphism such that $\phi' \circ \phi = 0$. Then $\phi(\Gamma) \subseteq \ker(\phi')$, but $\ker(\phi)$ is a convex subgroup by \cref{result:kernels-convex}, so $\Delta \subseteq \ker(\phi)$. Thus there exists a unique homomorphism $\bar{\phi}' : \Gamma'/\Delta \to \Gamma''$ such that the following diagram commutes.
\[ \begin{tikzcd} \Gamma \ar{r}{\phi} & \Gamma' \ar{r}{\pi} \ar{dr}[swap]{\phi'} & \Gamma'/\Delta \ar[dotted]{d}{\bar{\phi}'} \\ & & \Gamma'' \end{tikzcd} \]
In fact, $\bar{\phi}'$ is order preserving. We use the criterion of \cref{result:order-preserving-submonoid}. Suppose we have some $\pi(x) \in (\Gamma'/\Delta)_{\geq 0}$. Since $(\Gamma'/\Delta)_{\geq 0} = \pi(\Gamma_{\geq 0})$ by \cref{result:quotient-order}, we can assume that $x \geq 0$. Then $\bar{\phi}'(\pi(x)) = \phi'(x) \geq 0$ since $\phi'$ is order preserving, completing the proof. 
\end{proof}

\begin{remark}
It follows from \cref{result:categorical-zero-kernels-cokernels} that the category of totally ordered abelian groups has images and coimages, but it is clear from the proof of \cref{result:categorical-zero-kernels-cokernels} that the natural map from the coimage to the image will \emph{not} be an isomorphism in general. 
\end{remark}

\begin{lemma}
\label{result:monomorphism-injective}
Suppose $\phi : \Gamma \to \Gamma'$ is an order preserving homomorphism of totally ordered abelian groups. Then the following are equivalent.  
\begin{enumerate}[label=(\alph{*})]
\item $\phi$ is a monomorphism in the category of totally ordered abelian groups. 
\item $\phi$ is injective. 
\end{enumerate}
\end{lemma}

\begin{proof}
Injective maps are always monomorphisms in concrete categories, so we need only to prove that monomorphisms must be injective. But notice that $\phi$ coequalizes the inclusion $\ker(\phi) \hookrightarrow \Gamma$ and the zero map $\ker(\phi) \to \Gamma$, so if $\phi$ is a monomorphism, the inclusion must be the zero map, forcing $\ker(\phi) = 0$. 
\end{proof}

\begin{lemma}
Suppose $\phi : \Gamma \to \Gamma'$ is an order preserving homomorphism of totally ordered abelian groups. Then the following are equivalent.  
\begin{enumerate}[label=(\alph{*})]
\item $\phi$ is a normal monomorphism in the category of totally ordered abelian groups. 
\item $\phi$ is injective and $\phi(\Gamma)$ is a convex subgroup of $\Gamma'$. 
\end{enumerate}
\end{lemma}

\begin{proof}
If $\phi$ is injective and $\phi(\Gamma)$ is a convex subgroup of $\Gamma'$, then $\phi$ is a kernel of the projection $\pi : \Gamma' \to \Gamma'/\phi(\Gamma)$ and is therefore a regular monomorphism. Conversely, suppose $\phi$ is a normal monomorphism. Then $\phi$ is injective by \cref{result:monomorphism-injective}, and its image must be convex by \cref{result:kernels-convex}. 
\end{proof}

\begin{lemma}
	Suppose $\phi : \Gamma \to \Gamma'$ is an order preserving homomorphism of totally ordered abelian groups. Then the following are equivalent.  
	\begin{enumerate}[label=(\alph{*})]
		\item $\phi$ is an epimorphism in the category of totally ordered abelian groups. 
		\item The convex subgroup of $\Gamma'$ generated by $\phi(\Gamma)$ is $\Gamma'$. 
	\end{enumerate}
\end{lemma}

\begin{proof}
	Let $\Delta$ denote the convex subgroup of $\Gamma'$ generated by $\phi(\Gamma)$ (cf. \cref{result:intersection-convex}). Then $\phi$ equalizes the quotient map $\pi : \Gamma' \to \Gamma'/\Delta$ and the zero map $\Gamma' \to \Gamma'/\Delta$. If $\phi$ is an epimorphism, this means that the quotient map $\pi$ is the zero map, forcing $\Gamma'/\Delta = 0$, i.e., $\Delta = \Gamma'$. Conversely, suppose $\Delta = \Gamma'$ and that $\phi$ equalizes two order preserving homomorphisms of totally ordered abelian groups $\alpha, \beta : \Gamma' \to \Gamma''$. 
	\[ \begin{tikzcd} \Gamma \ar{r}{\phi} & \Gamma' \ar[bend left]{r}{\alpha} \ar[bend right]{r}[swap]{\beta} & \Gamma'' \end{tikzcd} \]
\end{proof}

% THE ABove is wrong

\begin{lemma}
	Suppose $\phi : \Gamma \to \Gamma'$ is an order preserving homomorphism of totally ordered abelian groups. Then the following are equivalent.  
	\begin{enumerate}[label=(\alph{*})]
		\item $\phi$ is a normal epimorphism in the category of totally ordered abelian groups. 
		\item $\phi$ is surjective.  
	\end{enumerate}
\end{lemma}

\begin{proof}
	If $\phi$ is surjective, then it is the cokernel of the inclusion $\ker(\phi) \hookrightarrow \Gamma$ and is therefore a normal epimorphism. Conversely, if $\phi$ is a normal epimorphism, it is a cokernel in the category of totally ordered abelian groups of some order preserving homomorphism $\alpha$ with codomain $\Gamma$. But this cokernel is precisely the quotient of $\Gamma$ by the convex subgroup generated by the set-theoretic image of $\alpha$, which is evidently surjective. 
\end{proof}

\begin{lemma}
Suppose $I$ is a index set of cardinality at least 2, and that $(\Gamma_i)_{i \in I}$ is a family of totally ordered abelian groups indexed by $I$. Then neither the product nor the coproduct of this family exists in the category of totally ordered abelian groups. 
\end{lemma}

\begin{proof}
Since $I$ has cardinality at least 2, it admits a nontrivial automorphism $\sigma : I \to I$.
\end{proof}

% INCOMPLETE

\section{Hahn embedding theorem}

\begin{proposition}[Hahn embedding theorem]
Let $\Gamma$ be a totally ordered abelian group. Then there exists an order embedding
\[ \begin{tikzcd} \Gamma \ar[hookrightarrow]{r} & \displaystyle \prod_{\Convex(\Gamma)}' \R. \end{tikzcd} \]
\end{proposition}

% INCOMPLETE


\printbibliography[heading=bibintoc,title={References}]

\end{document}